\documentclass[11pt]{article}
\usepackage[margin=1in]{geometry}
\usepackage{amsmath,amssymb,amsthm}
\usepackage{graphicx}
\usepackage{booktabs}
\usepackage{listings}
\usepackage{xcolor}
\usepackage[numbers]{natbib}
\usepackage{hyperref}

\lstset{basicstyle=\ttfamily\footnotesize,breaklines=true,columns=fullflexible,
  frame=single,extendedchars=true,
  literate={ℝ}{{$\mathbb{R}$}}1 {∀}{{$\forall$}}1 {→}{{$\to$}}1 {≤}{{$\le$}}1
    {↔}{{$\leftrightarrow$}}1 {∧}{{$\wedge$}}1 {·}{{$\cdot$}}1 {−}{{$-$}}1
    {–}{{--}}1 {≥}{{$\ge$}}1 {¬}{{$\neg$}}1 {≠}{{$\ne$}}1 {↦}{{$\mapsto$}}1
    {∈}{{$\in$}}1 {×}{{$\times$}}1 {`}{{\textasciigrave}}1}

\newtheorem{proposition}{Proposition}

\title{Persistence thresholds for a habitat-loss kernel: proposed formal statements for the Levins
threshold and power-law species--area bounds}
\author{Justin Hart\\Viridis LLC}
\date{27 November 2026}

\begin{document}
\maketitle

\begin{abstract}
Planners who weigh habitat loss usually want a threshold or a bound, not a forecast. Two textbook
models supply them. In the Levins metapopulation model with habitat loss, a population persists
only while the suitable fraction of habitat exceeds the ratio of extinction to colonisation rate.
Under a power-law species--area relationship $S=kA^z$ with $0<z\le1$, the fraction of species kept
is never below the fraction of area kept. This paper begins project P7 of Viridis Research by
writing both results as precise statements in Lean~4 over Mathlib, as candidates to be sealed as
machine-checkable targets. For each one we give the statement, every hypothesis, worked witnesses
and negative controls checked by Lean, and a list of what it does not settle. Both statements are
about real algebra only. The first concerns the sign of the Levins growth term, not solutions of
the differential equation. The second assumes the power law and says nothing about extinction
debt, which species are lost, or any value of $z$. We give no proof here; the proofs are later
tasks.
\end{abstract}

\section{Introduction}
How much habitat does a population need, and how much of a biota survives a given loss of area?
Two classical models give answers in the form of a threshold and a bound. The Levins
metapopulation model \citep{Levins1969}, extended to habitat destruction
\citep{Lande1987,NeeMay1992}, predicts that a species occupying a fraction $p$ of patches declines
whenever the suitable fraction $h$ falls to $e/c$ or below, where $e$ is the extinction rate and
$c$ the colonisation rate. The power-law species--area relationship \citep{Arrhenius1921} implies
that keeping a fraction $a$ of the area keeps the fraction $a^z$ of the species.

These are not new theorems. A habitat-loss kernel that cites them should still state exactly
which algebraic facts it relies on, under which hypotheses, and should have those facts checked
by machine. Project P7 seals these facts as formal statements. Tonight's task (P7-T001) is to
propose the two statements. Proving them is the job of P7-T002 and P7-T003.

\section{Background and related work}
\paragraph{Levins model with habitat loss.} With occupancy $p$, colonisation rate $c$, extinction
rate $e$ and suitable fraction $h$,
\[ \frac{dp}{dt} = c\,p\,(h-p) - e\,p = c\,p\,(p^*-p),\qquad p^* = h - e/c. \]
When habitat is intact ($h=1$) the equilibrium occupancy is $1-e/c$. Persistence requires
$h>e/c$, which is the same as saying the destroyed fraction $1-h$ is less than $1-e/c$. This is the
rule of thumb attributed to Lande and to Nee and May: a species is lost once the destroyed
fraction reaches the occupancy it had in intact habitat \citep{Lande1987,NeeMay1992}.

\paragraph{Species--area power law.} If $S=kA^z$, then $S(aA)/S(A)=a^z$. Both $k$ and $A$ cancel
\citep{Arrhenius1921}. For $0<z\le1$ and $0\le a\le1$ we have $a\le a^z\le1$, and the area fraction
needed to keep a fraction $s$ of species is $s^{1/z}\le s$.

\paragraph{Connectivity.} In the Levins model a corridor network can matter only through the
colonisation rate $c$. Project P6 of this programme treats route reliability. P7 makes no claim
about how any network changes $c$.

\section{Formal development}
Both modules live in their own namespaces, import only Mathlib, and work over $\mathbb{R}$. The
complete \texttt{Statement} definitions are given below exactly as proposed.

\subsection{P7.HabitatThreshold}
\begin{proposition}[proposed]
For all real $c,e,h$ with $c>0$, $e>0$: (1) $\mathrm{growth}(p)=c\,p\,(p^*-p)$ for every $p$;
(2) if $h\le e/c$ then $\mathrm{growth}(p)<0$ for every $p>0$; (3) if $e/c<h$ then $0<p^*<h$,
$\mathrm{growth}(p^*)=0$, $\mathrm{growth}>0$ on $(0,p^*)$ and $\mathrm{growth}<0$ on $(p^*,\infty)$;
(4) $e/c<h\iff 1-h<1-e/c$; (5) $e/c$ is strictly decreasing in $c$ and strictly increasing in $e$.
\end{proposition}
\begin{lstlisting}
namespace ViridisStatements.P7.HabitatThreshold

/-- Levins growth term with habitat loss: `c·p·(h − p) − e·p`. -/
def growth (c e h p : ℝ) : ℝ := c * p * (h - p) - e * p

/-- Nontrivial equilibrium occupancy `p* = h − e/c`. -/
noncomputable def pstar (c e h : ℝ) : ℝ := h - e / c

/-- **Proposed sealed target `P7.HabitatThreshold`.** -/
def Statement : Prop :=
  (∀ c e h : ℝ, 0 < c → 0 < e →
      ∀ p : ℝ, growth c e h p = c * p * (pstar c e h - p)) ∧
  (∀ c e h : ℝ, 0 < c → 0 < e → h ≤ e / c →
      ∀ p : ℝ, 0 < p → growth c e h p < 0) ∧
  (∀ c e h : ℝ, 0 < c → 0 < e → e / c < h →
      0 < pstar c e h ∧ pstar c e h < h ∧ growth c e h (pstar c e h) = 0 ∧
      (∀ p : ℝ, 0 < p → p < pstar c e h → 0 < growth c e h p) ∧
      (∀ p : ℝ, pstar c e h < p → growth c e h p < 0)) ∧
  (∀ c e h : ℝ, 0 < c → 0 < e → (e / c < h ↔ 1 - h < 1 - e / c)) ∧
  (∀ c e : ℝ, 0 < c → 0 < e →
      (∀ c' : ℝ, c < c' → e / c' < e / c) ∧ (∀ e' : ℝ, e < e' → e / c < e' / c))
\end{lstlisting}

\subsection{P7.SpeciesArea}
\begin{proposition}[proposed]
For all real $z$ with $0<z\le1$: (1) $a\le a^z\le1$ for $0\le a\le1$; (2) if $z<1$, $a<a^z$ for
$0<a<1$; (3) $a^z\le b^z$ for $0\le a\le b$; (4) for $0\le s\le1$, $(s^{1/z})^z=s$ and
$s^{1/z}\le s$. Powers are real powers (\texttt{Real.rpow}).
\end{proposition}
\begin{lstlisting}
namespace ViridisStatements.P7.SpeciesArea

/-- **Proposed sealed target `P7.SpeciesArea`.** Here `a ^ z` is the fraction of species kept
when a fraction `a` of the area is kept, and `s ^ (1 / z)` the area fraction needed to keep a
fraction `s` of the species. -/
def Statement : Prop :=
  (∀ z : ℝ, 0 < z → z ≤ 1 →
      ∀ a : ℝ, 0 ≤ a → a ≤ 1 → a ≤ a ^ z ∧ a ^ z ≤ 1) ∧
  (∀ z : ℝ, 0 < z → z ≤ 1 → z < 1 →
      ∀ a : ℝ, 0 < a → a < 1 → a < a ^ z) ∧
  (∀ z : ℝ, 0 < z → z ≤ 1 →
      ∀ a b : ℝ, 0 ≤ a → a ≤ b → a ^ z ≤ b ^ z) ∧
  (∀ z : ℝ, 0 < z → z ≤ 1 →
      ∀ s : ℝ, 0 ≤ s → s ≤ 1 → (s ^ (1 / z)) ^ z = s ∧ s ^ (1 / z) ≤ s)
\end{lstlisting}

\section{Methods: checks before proposing}
We checked each step from the task brief by hand. Since $c\neq0$, $c\,p\,(h-e/c-p)=c\,p\,(h-p)-e\,p$.
Below the threshold, $p^*\le0<p$ gives $p^*-p<0$ while $c\,p>0$. Above it, $p>p^*>0$ gives $c\,p>0$.
For the power law, $x\mapsto a^x$ is antitone when $0<a\le1$, and $1/z\ge1$ when $0<z\le1$.

To avoid proposing a false target, we also wrote a private Lean proof of each \texttt{Statement}
text in a scratch file. That file is not part of this package and supports no claim here. It
turned up one detail worth recording. In the pinned Mathlib,
\texttt{Real.rpow\_le\_rpow\_of\_exponent\_ge} needs a strictly positive base, so the case $a=0$
of part~1 of \texttt{P7.SpeciesArea} needs the primed variant
\texttt{Real.rpow\_le\_rpow\_of\_exponent\_ge'}. We also confirmed that none of
\texttt{norm\_num}, \texttt{simp}, \texttt{aesop}, \texttt{decide}, \texttt{grind},
\texttt{positivity} or \texttt{linarith} closes either statement after unfolding, so each one is
a non-trivial target.

\section{Results: sanity checks and controls}
Each module contains \texttt{example}s that Lean checks when the module compiles
(Table~\ref{tab:checks}). \texttt{decide} cannot evaluate real arithmetic or real powers, so the
checks use \texttt{norm\_num}, with \texttt{Real.rpow\_mul} and \texttt{Real.rpow\_natCast} for the
powers.

\begin{table}[h]\centering
\begin{tabular}{@{}lll@{}}\toprule
Module & Check & Value \\\midrule
HabitatThreshold & $c=2,e=1,h=1$: $p^*$ & $1/2$ \\
 & growth at $p=1/4$, $h=1$ & $1/8>0$ \\
 & growth at $p=1/4$, $h=1/2$ & $-1/8<0$ \\
 & control: growth at $p=0$, $h=1/2$ & $0$, not $<0$ \\
 & control: $e=0$, $c=2\to3$ & $0/3=0/2$ \\
SpeciesArea & $z=1/2$, $a=1/4$: $a^z$ & $1/2$ \\
 & $z=1/2$, $s=1/2$: $s^{1/z}$ & $1/4$ \\
 & control: $z=2$, $a=1/2$: $a^z$ & $1/4<1/2$ \\\bottomrule
\end{tabular}
\caption{Witnesses and negative controls included as checked examples.}\label{tab:checks}
\end{table}

The negative controls show that the hypotheses are needed. Without $0<p$, part~2 of
\texttt{P7.HabitatThreshold} fails at $p=0$. Without $z\le1$, part~1 of \texttt{P7.SpeciesArea}
fails at $z=2$, $a=1/2$.

\section{Limitations}
A future proof of these statements would establish the algebra of two models and nothing about
nature. \texttt{P7.HabitatThreshold} concerns the sign of the growth term, not solutions of
$dp/dt=\mathrm{growth}$, and $c$, $e$ and $h$ are field inputs whose values it does not
constrain. The Levins model assumes identical patches, mean-field colonisation, constant rates and
no stochastic extinction. \texttt{P7.SpeciesArea} assumes $S=kA^z$, an empirical regularity. It says
nothing about extinction debt, which species are lost, where the remaining habitat lies, or any
value of $z$.

On the programme's frontier duty for premise \texttt{M-Nucleation}, we did not register a test
plan (see the notebook). The public data we found from an applied-nucleation experiment in Costa
Rica \citep{Zahawi2013} records recruits per tree island. It does not record whether planted
clusters below a critical size shrink, so it cannot decide the nucleation premise.

\section{Reproducibility}
The proposals are \texttt{proposals/P7.HabitatThreshold.lean} and
\texttt{proposals/P7.SpeciesArea.lean}. They compile with \texttt{lean} under
\texttt{leanprover/lean4:v4.28.0} and Mathlib commit
\texttt{8f9d9cff6bd728b17a24e163c9402775d9e6a365}, and the examples in them are the checks in
Table~\ref{tab:checks}. Rationales are in the matching \texttt{.md} files. No data or code beyond
these files was used.

\section{Changelog}
\begin{itemize}
\item 2026-11-27 (P7-T001): first version. Proposed \texttt{P7.HabitatThreshold} and
\texttt{P7.SpeciesArea}, with sanity checks, negative controls and rationales. No proofs.
\end{itemize}

\bibliographystyle{plainnat}
\bibliography{refs}
\end{document}
